% =====================================================================
% 02_methodology.tex  --  Detailed methodology (Exp. 4)
% =====================================================================

\section{Methodology}
\label{sec:dproxy-method}

The analysis has two parts that use the same 500 images but not the same decode.
Part~one (Section~\ref{subsec:dproxy-danalysis}) reads the contrastive adjustment
$d = E - A$ from logits stored along the real VCD or SID decoding trajectory (the
path chosen by $C = 2E - A$). Part~two (Section~\ref{subsec:dproxy-proxy})
performs an independent expert-only decode with random noise added to the logits,
and asks whether that noise stand-in reproduces the real methods' effect on
hallucination. Only the image set is shared; the two parts are separate decoding
runs. We first fix the common setup.

\subsection{Shared setup}
\label{subsec:dproxy-setup}

\textbf{Models and data.} We use LLaVA-1.5-7B and Qwen2.5-VL-7B-Instruct, both in
\texttt{float16} with \texttt{sdpa} attention, on a fixed set of 500 MSCOCO~\citep{lin2014coco}
\texttt{val2017} images (the same images and order for every run, so all
comparisons are paired). The prompt is ``Describe this image in detail'' and
decoding is greedy with \texttt{max\_new\_tokens} $=256$.

\textbf{Contrastive branches.} The expert branch $E$ scores the next token from
the clean image. The amateur branch $A$ scores it from a degraded image. For VCD
the degradation is forward diffusion noise at step 500; for SID it is attention
pruning that keeps a random $12.5\%$ of image tokens from layer two onward. The
contrastive score is $C = 2E - A$ (that is $\alpha = 1$), followed by the
adaptive plausibility constraint (APC), which removes any token whose expert
logit is below $\log\beta + \max_w E_t(w)$ with $\beta = 0.2$. These settings
match the reference implementations and the other experiments in this paper.

\textbf{What is recorded.} At every decoding step we store the top-30 token ids
and logits of both the expert and the amateur branch, together with the token
the method finally emits. All object labelling is done afterward from the
generated caption text, never from raw sub-word tokens, which avoids the sub-word
mislabelling problem (for example the ``ship'' inside ``companionship'' or a
``bear'' split off from ``teddy bear'').

\textbf{Object labelling.} We use the standard CHAIR
protocol~\citep{rohrbach2018chair}. Each image has a
ground-truth object set built from its COCO instance-segmentation labels plus the
objects named in its five human reference captions, mapped to the 80 MSCOCO
categories through the CHAIR synonym list. An object the model mentions is
\emph{real} if its category is in that set and \emph{hallucinated} if it is not.

\subsection{Part one: the contrastive adjustment \texorpdfstring{$d = E - A$}{d = E - A}}
\label{subsec:dproxy-danalysis}

For a given token we read its expert logit $E$ and amateur logit $A$ from the
stored top-30 lists and form $d = E - A$. Because $C = E + d$, a positive $d$
means the token is amplified by contrastive decoding and a negative $d$ means it
is suppressed. We compute $d$ over three populations of object tokens, from the
cleanest to the broadest, so that the conclusion does not rest on one arbitrary
cut-off.

\begin{enumerate}
  \item \textbf{Emitted object words (primary).} These are the object words the
  model actually wrote, labelled at the caption level, so every real or
  hallucinated tag is exact. For each such word we take the step that produced
  its first token and read $E$ and $A$ for that token.
  \item \textbf{Top-10 expert candidates (robustness).} At every step we take the
  ten highest tokens by expert logit, keep those that are word-initial COCO
  object words, and label each by ground-truth membership. This is not
  conditioned on a word winning, so it removes any selection effect of the
  emitted set.
  \item \textbf{Top-30 expert candidates (robustness).} The same as above with a
  wider window, to check that the picture is stable as the candidate set grows.
\end{enumerate}

For each population we report, separately for real and hallucinated objects, the
mean of $d$ and the fraction of tokens with $d > 0$ (the share that contrastive
decoding amplifies rather than suppresses). We visualise the full spread with
per-token bar plots, one bar per token, coloured red where $d > 0$ and blue where
$d < 0$, with the group mean drawn as a dashed line.

\textbf{Two honest notes.} First, if a token is not present in both the expert and
the amateur stored top-30 lists, its $d$ cannot be formed and the token is
skipped; this is a small fraction and is applied identically to both branches, so
it cannot bias the real-versus-hallucinated comparison. Second, the candidate
populations use single-token labels, so a word-initial ``bear'' that begins
``teddy bear'' would be read as the animal; this affects well under one percent
of candidates and hits both branches equally. The emitted population has no such
caveat, which is why we treat it as primary.

\subsection{Part two: the noise proxy}
\label{subsec:dproxy-proxy}

The measurements in Section~\ref{subsec:dproxy-danalysis} suggest a specific
reading of the contrastive term. The adjustment $d = E - A$ that contrastive
decoding adds to the expert logits does not behave like a targeted,
content-aware signal: it does not single out hallucinated objects for
suppression, and across tokens it resembles a small, roughly zero-centred
perturbation layered on top of the expert scores. This motivates a direct test.
If the amateur branch contributes, in effect, only a perturbation of a certain
size, then adding a matched amount of \emph{random} noise to the same expert
logits should produce a similar outcome, and the amateur branch would not be
doing anything a generic perturbation could not.

We therefore build a deliberately minimal control, the noise proxy. We discard
the amateur branch entirely. Starting from the plain expert logits $E$, we add
independent random noise to the object-category logits, with magnitude matched to
the measured distribution of $d$ for the real method, and then apply the same
adaptive plausibility constraint and greedy decoding. The proxy has no access to
the degraded image and performs no contrast; its only ingredients are the expert
logits and random noise of the same size that contrastive decoding actually adds.
We evaluate it on the standard CHAIR benchmark under identical conditions and
compare against the real VCD and SID runs. The logic is straightforward: if
contrastive decoding worked through a genuine image-grounded mechanism, matched
random noise should not reproduce its effect on hallucination; if the proxy lands
at the same CHAIR scores as the real methods, the measured effect of contrastive
decoding is statistically indistinguishable from adding random noise of the same
size, and the amateur branch is not contributing targeted, hallucination-reducing
information.

\textbf{Object-token vocabulary.} We first collect the set of vocabulary token
ids that are word-initial and spell an MSCOCO object word (348 tokens for
Qwen; a comparable set for LLaVA). These are the only tokens the proxy perturbs,
because these are the tokens whose logits determine CHAIR. Restricting the noise
to object tokens is deliberate and conservative: it gives the proxy the best
possible chance to mimic a hallucination-relevant intervention.

\textbf{Which $d$ distribution the noise is matched to.} The noise parameters in
Table~\ref{tab:dproxy-noise}, and the empirical values that Proxy~B resamples, are
computed over a single broad pool: every object-category token that appears
anywhere in the stored top-30 expert candidate list, at every step, with real and
hallucinated tokens combined (about $114{,}000$ tokens per model). This pool is
deliberately broader than the top-10, top-30, or emitted populations used in the
$d$-analysis of Section~\ref{subsec:dproxy-danalysis}, which is why its mean does
not equal any single value in those tables. Using one pooled distribution keeps
the proxy agnostic to whether a token is real or hallucinated, which is the point
of the control.

\textbf{Proxy A (Gaussian, magnitude-matched).} We run expert-only greedy
decoding. At each step we add to every object-token logit an independent draw
from a Gaussian whose mean and standard deviation equal the measured pooled mean
and standard deviation of $d$ for the corresponding real method
(Table~\ref{tab:dproxy-noise}). We then apply the same APC and take the argmax.
There is no amateur branch and no second image pass.

\textbf{Proxy B (empirical bootstrap).} Identical to Proxy A, except the bonus
for each object token is drawn by resampling with replacement from the empirical
set of measured $d$ values for that method. This reproduces the real
distribution of $d$, including its skew, not just its first two moments.

\textbf{Matching the magnitude.} Table~\ref{tab:dproxy-noise} lists the noise
parameters. The sign of $d$ differs by model: on LLaVA the mean of $d$ is
positive, so the proxy on average raises object logits; on Qwen it is negative,
so the proxy on average lowers them. In both cases the proxy matches the real
method's own measured magnitude, so it is a like-for-like stand-in.

\begin{table}[t]
\centering
\small
\caption{Noise parameters for the proxies, taken from the measured distribution
of $d = E - A$ over object tokens. Proxy~A draws from a Gaussian with these mean
and standard deviation; Proxy~B resamples the empirical $d$ values (same mean and
standard deviation by construction).}
\label{tab:dproxy-noise}
\begin{tabular}{llcc}
\toprule
Model & Method & mean of $d$ & std of $d$ \\
\midrule
LLaVA-1.5-7B & VCD & $+0.199$ & $0.853$ \\
 & SID & $+0.300$ & $0.903$ \\
Qwen2.5-VL-7B & VCD & $-0.846$ & $1.531$ \\
 & SID & $-0.663$ & $1.558$ \\
\bottomrule
\end{tabular}
\end{table}

\textbf{Scoring.} Every proxy run produces 500 captions on the same images, which
we score with the standard CHAIR metric using the identical scorer applied to the
real greedy, VCD, and SID runs. We report CHAIR-S (the percentage of captions
containing at least one hallucinated object) and CHAIR-I (the percentage of
mentioned object instances that are hallucinated). Lower is better for both.

\textbf{Reproducibility note.} No global random seed was fixed. The noise draws
in the proxies and the stochastic degradations in the real methods therefore
differ across runs, but every quantity is aggregated over 500 images and tens of
thousands of tokens, a regime where this variation is negligible relative to the
reported effects.
